% ============================================================
%  MODULE IV — Fundamental Analysis
% ============================================================

\chapter{Fundamental Analysis}

\begin{notebox}
Fundamental analysis asks: \textbf{what is this asset actually worth?}
It grounds investment decisions in business reality rather than price patterns alone.
Warren Buffett, Charlie Munger, and most long-term investors use this approach.
\end{notebox}

% ────────────────────────────────────────────────────────────
\section{The Three Financial Statements}

Every listed company publishes three core statements. Learn to read all three.

\subsection{Income Statement (Profit \& Loss)}
Shows \textbf{revenues, expenses, and profit} over a period (quarterly or annually).

\begin{center}
\begin{tabular}{>{\bfseries}p{5cm} p{7.5cm}}
\toprule
Line Item & What It Means \\
\midrule
Revenue / Sales & Total money earned from operations \\
Cost of Goods Sold (COGS) & Direct costs of producing goods/services \\
Gross Profit & Revenue $-$ COGS \\
Operating Expenses (OpEx) & R\&D, SG\&A, depreciation \\
EBIT & Earnings Before Interest and Tax \\
EBITDA & EBIT + Depreciation + Amortisation; proxy for operating cash \\
Net Income & Final profit after interest and tax \\
EPS & Earnings Per Share = Net Income $\div$ Shares Outstanding \\
\bottomrule
\end{tabular}
\end{center}

\subsection{Balance Sheet}
A \textbf{snapshot} of what the company owns (assets) and owes (liabilities) at one moment.

\begin{formulabox}
\[
  \text{Assets} = \text{Liabilities} + \text{Shareholders' Equity}
\]
This equation always holds. Equity is the residual claim.
\end{formulabox}

\subsection{Cash Flow Statement}
Tracks \textbf{actual cash movement}, split into three activities:
\begin{itemize}
  \item \textbf{Operating Cash Flow (OCF)} — cash from core business. Most important signal of health.
  \item \textbf{Investing Cash Flow} — capex, acquisitions, asset sales.
  \item \textbf{Financing Cash Flow} — debt issuance/repayment, dividends, share buybacks.
\end{itemize}

\begin{warningbox}
A company can show positive net income while \emph{burning cash}. Always check the
cash flow statement. Enron's fraud persisted partly because investors ignored cash flow signals.
\end{warningbox}

% ────────────────────────────────────────────────────────────
\section{Key Financial Ratios}

\subsection{Valuation Ratios}

\begin{center}
\begin{tabular}{>{\bfseries}p{3cm} p{5cm} p{5cm}}
\toprule
Ratio & Formula & What It Tells You \\
\midrule
P/E & Price $\div$ EPS & How much investors pay per \$1 of earnings \\
P/S & Market Cap $\div$ Revenue & Useful for unprofitable companies \\
P/B & Price $\div$ Book Value per Share & Premium over net assets \\
EV/EBITDA & Enterprise Value $\div$ EBITDA & Capital-structure-neutral valuation \\
PEG & P/E $\div$ EPS Growth Rate & Adjusts P/E for growth \\
\bottomrule
\end{tabular}
\end{center}

\begin{formulabox}
\textbf{Enterprise Value (EV):}
\[
  EV = \text{Market Cap} + \text{Total Debt} - \text{Cash \& Equivalents}
\]
EV represents the \emph{total cost to acquire} a business.
\end{formulabox}

\subsection{Profitability Ratios}

\begin{center}
\begin{tabular}{>{\bfseries}p{3.5cm} p{5.5cm} p{4.5cm}}
\toprule
Ratio & Formula & Interpretation \\
\midrule
Gross Margin & Gross Profit $\div$ Revenue & Pricing power, cost efficiency \\
Operating Margin & EBIT $\div$ Revenue & Core business profitability \\
Net Margin & Net Income $\div$ Revenue & Bottom-line efficiency \\
ROE & Net Income $\div$ Equity & Return on shareholders' capital \\
ROA & Net Income $\div$ Total Assets & Asset efficiency \\
ROIC & NOPAT $\div$ Invested Capital & True economic return on capital deployed \\
\bottomrule
\end{tabular}
\end{center}

\subsection{Liquidity and Leverage Ratios}

\begin{center}
\begin{tabular}{>{\bfseries}p{3.5cm} p{5.5cm} p{4.5cm}}
\toprule
Ratio & Formula & Interpretation \\
\midrule
Current Ratio & Current Assets $\div$ Current Liabilities & Short-term solvency \\
Quick Ratio & (Current Assets $-$ Inventory) $\div$ CL & Stricter solvency test \\
Debt-to-Equity & Total Debt $\div$ Equity & Financial leverage \\
Net Debt/EBITDA & Net Debt $\div$ EBITDA & Years to repay debt from operations \\
Interest Coverage & EBIT $\div$ Interest Expense & Ability to service debt \\
\bottomrule
\end{tabular}
\end{center}

% ────────────────────────────────────────────────────────────
\section{Discounted Cash Flow (DCF) Valuation}

The theoretical gold standard: value a company as the present value of all future free
cash flows (FCF), discounted at the \textbf{Weighted Average Cost of Capital (WACC)}.

\begin{formulabox}
\[
  \text{Intrinsic Value} = \sum_{t=1}^{n} \frac{FCF_t}{(1 + WACC)^t} + \frac{TV}{(1 + WACC)^n}
\]
where $TV$ = terminal value (value of all cash flows beyond the forecast period).

\textbf{Free Cash Flow:}
\[
  FCF = \text{Operating Cash Flow} - \text{Capital Expenditure}
\]

\textbf{WACC:}
\[
  WACC = \frac{E}{V} \cdot r_e + \frac{D}{V} \cdot r_d \cdot (1 - T)
\]
where $E/V$ = equity weight, $D/V$ = debt weight, $r_e$ = cost of equity,
$r_d$ = cost of debt, $T$ = tax rate.
\end{formulabox}

\begin{notebox}
DCF is highly sensitive to assumptions — a 1\% change in WACC or growth rate can
change the output by 30--50\%. Use it as a \emph{range}, not a precise number.
``Garbage in, garbage out.''
\end{notebox}

% ────────────────────────────────────────────────────────────
\section{Qualitative Factors}

Numbers are necessary but not sufficient. Also assess:

\begin{itemize}
  \item \textbf{Competitive moat} — durable competitive advantage (brand, network effects,
        switching costs, cost advantages, intangible assets). Popularised by Warren Buffett.
  \item \textbf{Management quality} — capital allocation history, insider ownership, communication.
  \item \textbf{Industry dynamics} — growth trajectory, competitive intensity (Porter's Five Forces).
  \item \textbf{Regulatory risk} — antitrust, environmental, financial regulation exposure.
  \item \textbf{Balance sheet resilience} — can the company survive a recession or industry shock?
\end{itemize}

% ────────────────────────────────────────────────────────────
\section{Reading Earnings Reports}

Every quarter, public companies release:
\begin{enumerate}
  \item \textbf{Earnings Release} — headline EPS, revenue, and guidance.
  \item \textbf{10-Q / 10-K} (US) — detailed SEC filings with full financial statements.
  \item \textbf{Earnings Call} — management discusses results and takes analyst questions.
\end{enumerate}

\begin{notebox}
Markets react to \textbf{surprises vs. expectations}, not absolute numbers. A company with
great results can fall in price if the results missed consensus estimates. Track
\textbf{EPS surprise} and \textbf{revenue surprise} alongside guidance changes.
\end{notebox}
